\section{\personNovikov{} \personFloer{}链与\termNovikovRing{}}
\label{sec:novikov-floer-chains-ring}

设 $\phi=\phi_H^1\in\operatorname{Ham}(M,\omega)$ 非退化。对每个这样的
$H:S^1\times M\to\mathbb R$，我们知道 $\operatorname{Per}(H)$ 的基数有限。
考虑由 $\mathcal A_H$ 的临界点集
\[
 \operatorname{Crit}\mathcal A_H
 =\{[z,w]\in\widetilde{\mathcal L}_0(M)\mid z\in\operatorname{Per}(H)\}
\]
生成的自由 $\mathbb Q$-向量空间。为了正确定义\personFloer{}边界算子，需要关于元素
$[z,w]$ 的作用量 $\mathcal A_H([z,w])$ 所给出的实滤过，向下完备化这个向量空间。

\begin{definition}
\label{def:novikov-floer-chain}
考虑形式和
\begin{equation}
 \beta=\sum_{[z,w]\in\operatorname{Crit}\mathcal A_H}
 a_{[z,w]}[z,w],\qquad a_{[z,w]}\in\mathbb Q.
\label{eq:19-formal-floer-sum}
\end{equation}
\begin{enumerate}[label=(\arabic*)]
\stfhSubIndex{I142}{NuoweikefuFuluoerlian}{诺维科夫弗洛尔链（Novikov Floer chain）}{Zhiji}{支集（support）}
\item 定义 $\beta$ 的\emph{支集}为
\[
 \operatorname{supp}(\beta):=
 \{[z,w]\in\operatorname{Crit}\mathcal A_H\mid
 a_{[z,w]}\neq0\text{，在和式~\eqref{eq:19-formal-floer-sum} 中}\}.
\]
\stfhSubIndex{I140}{NuoweikefuFuluoerlian}{诺维科夫弗洛尔链（Novikov Floer chain）}{Shengchengyuan}{生成元（generator）}
\item 称满足 $a_{[z,w]}\neq0$ 的 $[z,w]$ 为和式 $\beta$ 的\emph{生成元}，并记作
\[
 [z,w]\in\operatorname{supp}\beta.
\]
在这种情形下，也称 $[z,w]$ 对 $\beta$ \emph{有贡献}。
\item 若形式和 $\beta$ 对任意 $\lambda\in\mathbb R$ 都满足
\begin{equation}
 \#\bigl(\operatorname{supp}(\beta)\cap
 \{[z,w]\mid\mathcal A_H([z,w])\geq\lambda\}\bigr)<\infty,
\label{eq:19-novikov-finiteness}
\end{equation}
\stfhIndex{I139}{NuoweikefuFuluoerlian}{诺维科夫弗洛尔链（Novikov Floer chain）}
则称 $\beta$ 为\emph{\personNovikov{} \personFloer{}链}（简称\emph{\personFloer{}链}）。称此条件为\termNovikovFinitenessCondition{}，并以 $CF_*(H)$ 表示\personFloer{}链的集合。
\end{enumerate}
\end{definition}

注意，按此定义，$CF_*(H)$ 是一个 $\mathbb Q$-向量空间，一般总是无限维的；只有当
$(M,\omega)$ 辛非球面时，它才是有限维的。给每个生成元
$[z,w]\in\operatorname{Crit}\mathcal A_H$ 配上\personConley{}--\personZehnder{}指标
$\mu_H([z,w])$，便赋予 $CF_*(H)$ 分次向量空间结构。

现在考虑\personFloer{}链
\[
 \beta=\sum a_{[z,w]}[z,w],\qquad a_{[z,w]}\in\mathbb Q.
\]
下面的概念是借助\personFloer{}同调建立极小极大临界点理论的关键。

\stfhSubIndex{I141}{NuoweikefuFuluoerlian}{诺维科夫弗洛尔链（Novikov Floer chain）}{Feng}{峰（peak）}
\begin{definition}
\label{def:floer-chain-level-peak}
设 $0\neq\beta\in CF_*(H)$ 是\personFloer{}链。定义链 $\beta$ 的\emph{层级}为
\[
 \lambda_H(\beta)=\max_{[z,w]}
 \{\mathcal A_H([z,w])\mid[z,w]\in\operatorname{supp}(\beta)\},
\]
并令 $\lambda_H(0)=-\infty$。若生成元 $[z,w]\in\beta$ 满足
$\mathcal A_H([z,w])=\lambda_H(\beta)$，则称它为 $\beta$ 的一个\emph{峰}，并记作
$\operatorname{peak}(\beta)$。
\end{definition}

我们强调，正是定义 19.1.1 中的\personNovikov{}有限性条件~\eqref{eq:19-novikov-finiteness}
保证了 $\lambda_H(\beta)$ 良定义。下述引理说明\personNovikov{}覆盖空间的定义是最优的。

\begin{lemma}
\label{lem:unique-peak-fixed-orbit}
设 $\beta\neq0$ 是齐次\personFloer{}链。则在一个固定周期轨道上，$\beta$ 的峰是唯一的。
\end{lemma}

\begin{proof}
设 $[z,w]$ 与 $[z,w']$ 是 $\beta$ 的两个这样的峰。于是
\[
 \mathcal A_H([z,w])=\lambda_H(\beta)=\mathcal A_H([z,w']),
\]
从而 $\omega([w])=\omega([w'])$。由齐次性假设还有
\[
 \mu_H([z,w])=\mu_H([z,w']).
\]
根据 $\Gamma$-等价类的定义，$[z,w]=[z,w']$，证明完毕。
\end{proof}

到目前为止，我们已把 $CF_*(H)$ 定义为以
$\operatorname{Crit}\mathcal A_H$ 为生成集的 $\mathbb Z$-分次
$\mathbb Q$-向量空间。现在说明如何像 (Fl89b)、(HS95) 那样，把 $CF(H)$ 描述成\emph{\personNovikov{}环}上的模。

回忆 $\Gamma$ 是交换群。考虑由有限和
\[
 R=\sum_{i=1}^{k}r_iq^{A_i}\in\mathbb Q[\Gamma]
\]
组成的群环 $\mathbb Q[\Gamma]$，并定义其支集为
\[
 \operatorname{supp}R=
 \{A\in\Gamma\mid A=A_i\text{，且在此和式中 }r_i\neq0\}.
\]
回忆赋值 $\nu:\Gamma\to\mathbb R$ 与次数映射 $d:\Gamma\to\mathbb Z$
\footnote{译者注：原书此处把次数映射的陪域印作 $\mathbb R$；但此前式 (18.2.13)
已定义 $d:\Gamma\to\mathbb Z$、$d(A)=2c_1(A)$，译文依该定义改正。}
。
赋值 $\nu:\mathbb Q[\Gamma]\to\mathbb R$ 定义为
\[
 \nu(R)=\nu\left(\sum_{i=1}^{k}r_iq^{A_i}\right)
 :=\min\{\omega(A)\mid A\in\operatorname{supp}R\}.
\]
它满足如下非阿基米德三角不等式：
\begin{equation}
 \nu(R_1+R_2)\geq\min\{\nu(R_1),\nu(R_2)\},
\label{eq:19-nonarch-triangle}
\end{equation}
因而在 $\mathbb Q[\Gamma]$ 上诱导由下列距离给出的自然度量拓扑
\footnote{译者注：原书在此写 $e^{\nu(R_1-R_2)}$。若赋值趋于 $+\infty$ 表示\personNovikov{}级数收敛，则通常相应距离应含负指数；原式与紧接着按此赋值作“向上完备化”
的约定表面上不相容。信息不足以唯一判定作者的符号规范，故正文照录。}
：
\[
 d(R_1,R_2):=e^{\nu(R_1-R_2)}.
\]

\begin{definition}
\label{def:novikov-ring-upward-completion}
\personNovikov{}环是 $\mathbb Q[\Gamma]$ 关于赋值
$\nu:\mathbb Q[\Gamma]\to\mathbb R$ 的\emph{向上}完备化
$\mathbb Q[\Gamma\rangle$，记作 $\Lambda_\omega=\Lambda_\omega^\uparrow$。
\end{definition}

更具体地，
\[
 \Lambda_\omega=
 \left\{\sum_{A\in\Gamma}r_Aq^A\ \middle|\
 \forall\lambda\in\mathbb R,
 \#\{A\in\Gamma\mid r_A\neq0,\ \omega(A)<\lambda\}<\infty\right\}.
\]
在 $\operatorname{Crit}\mathcal A_H$ 上定义“粘合一个球面”的 $\Gamma$-作用
\[
 [z,w]\longmapsto[z,w\#(-A)],
\]
进而以和 $\Lambda_\omega$ 的卷积积在 $CF(H)$ 上诱导 $\Lambda_\omega$-模结构：
\begin{equation}
\begin{aligned}
 &\left(\sum_{A\in\Gamma}r_Aq^A\right)
 \cdot\left(\sum a_{[z,w]}[z,w]\right)\\
 &\quad:=\sum_{[z,w']}
 \left(\sum_{\substack{A,[z,w]\,;\\{}[z,(-A)\#w]=[z,w']}}
 r_Aa_{[z,(-A)\#w]}\right)[z,w'].
\end{aligned}
\label{eq:19-convolution-action}
\end{equation}
按照这个定义，$q^A[z,w]=[z,(-A)\#w]$。换言之，因为这里讨论的是同调而非上同调，
我们改变了符号。

以下我们力求始终以 $CF(H)$ 表示 $\Lambda_\omega$-模，以 $CF_*(H)$ 表示分次
$\mathbb Q$-向量空间。作用量泛函在 $CF_*(H)$ 上给出自然滤过：给定任意
$\lambda\in\mathbb R\setminus\operatorname{Spec}(H)$，定义
\[
 CF_*^\lambda(H)=\{\alpha\in CF_*(H)\mid\lambda_H(\alpha)\leq\lambda\},
\]
并以
\[
 i_\lambda:CF_*^\lambda(H)\longrightarrow CF_*(H)
\]
表示自然包含同态。

稍后会看到，把 $q^A=T^{\omega(A)}e^{c_1(A)}$ 中的形式变量 $T$ 与 $e$ 分开有时
很有用。记 $\Lambda_\omega^{(0)}$ 为 $\Lambda_\omega$ 的零次部分，即满足\personNovikov{}有限性条件的形式级数
\[
 \sum_{\substack{A\in\Gamma\\c_1(A)=0}}a_Aq^A
 =\sum_{\substack{A\in\Gamma\\c_1(A)=0}}a_AT^{\omega(A)}
\]
的集合。显然，我们需要验证下述结论。

\begin{proposition}
\label{prop:degree-zero-novikov-field}
$\Lambda_\omega^{(0)}$ 是域，并且
\[
 \Lambda_\omega=\Lambda_\omega^{(0)}[e,e^{-1}],
\]
即 $\Lambda_\omega$ 是域 $\Lambda_\omega^{(0)}$ 上关于 $e,e^{-1}$ 的多项式环。
\end{proposition}

\begin{exercise}
\label{ex:prove-degree-zero-novikov-field}
证明此命题。
\end{exercise}
